\documentclass[11pt]{article}

\usepackage[margin=0.92in]{geometry}
\usepackage{amsmath,amssymb,amsthm,mathtools}
\usepackage{graphicx}
\usepackage[hidelinks]{hyperref}

\newtheorem{theorem}{Theorem}
\newtheorem{lemma}{Lemma}
\newtheorem{remark}{Remark}
\newcommand{\N}{\mathbb N}
\newcommand{\R}{\mathbb R}
\newcommand{\uu}[1]{(\cos #1,\sin #1)}

\title{No Finite Window Forces Banach-Valued Fekete Convergence}
\author{Automated Math Discovery Candidate Counterexample}
\date{August 11, 2026}

\begin{document}
\maketitle

\begin{center}
\fbox{\parbox{0.92\textwidth}{\textbf{Status.}
Candidate full negative answer to Question 3.2 of Kulikov--Shao,
arXiv:2411.17380v1, likely valid pending expert review.  For every proposed
finite window function, the required inequalities admit a nonconvergent
sequence already in the Euclidean plane.}}
\end{center}

\section*{Source Question}

Kulikov and Shao prove a Banach-valued Fekete lemma when
\(\|v_{n+m}\|\leq\|v_n+v_m\|\) is imposed for all pairs \cite{KS}.  After
constructing planar counterexamples for several wide but prescribed windows,
they ask on page 8:

\begin{center}
\includegraphics[width=\textwidth]{figures/open_problem_crop.png}
\end{center}

In other words, for a fixed uniformly convex space \(X\), can any function
\(f:\N\to\N\) make the restricted inequalities
\[
 \|v_{n+m}\|\leq\|v_n+v_m\| \qquad(n\leq m\leq f(n))
\]
force convergence of \(v_n/n\)?  The next theorem answers no for
\(X=\R^2\), hence answers the question negatively.

\section*{Answer and Proof Intuition}

\begin{theorem}[No finite window in the Euclidean plane]\label{thm:main}
For every function \(f:\N\to\N\), there is a sequence
\((v_n)_{n\geq1}\subset\R^2\) such that
\[
 \|v_{n+m}\|_2\leq\|v_n+v_m\|_2
 \quad\text{whenever }n\leq m\leq f(n),
\]
but \((v_n/n)\) does not converge.
\end{theorem}

The construction uses two scales.  First, partition \(\N\) into finite blocks
chosen recursively from \(f\), so that an allowed pair can meet at most two
consecutive blocks.  Second, rotate the direction by \(c/k\) between blocks
and add radial slack proportional to the tail
\(\sum_{j\geq k}j^{-2}\).  The angular increments have divergent sum, so the
directions rotate forever; their squares have convergent sum, so a vanishing
amount of radial slack can pay for the Euclidean cosine loss.

\section*{Proof}

Fix \(f:\N\to\N\).  Put \(N_1=1\), \(N_2=2\), and recursively choose
\(N_{k+2}\) so large that
\begin{equation}\label{eq:boundary}
 N_{k+2}>N_{k+1}
 \quad\text{and}\quad
 N_{k+2}>\max\{f(n):1\leq n<N_{k+1}\}.
\end{equation}
This is possible because the set in the maximum is finite.  Let
\[
 B_k=\{n\in\N:N_k\leq n<N_{k+1}\}.
\]
Every block is nonempty.  More importantly, if \(n\in B_k\) and
\(n\leq m\leq f(n)\), then \eqref{eq:boundary} gives
\begin{equation}\label{eq:locality}
                  m<N_{k+2},\qquad m\in B_k\cup B_{k+1}.
\end{equation}

Set \(c=1/4\), and for \(k\geq1\) define
\begin{equation}\label{eq:parameters}
 \alpha_1=0,\qquad
 \alpha_{k+1}=\alpha_k+\frac{c}{k},
 \qquad
 s_k=2c^2\sum_{j=k}^{\infty}\frac1{j^2}.
\end{equation}
Then \(s_k\downarrow0\),
\begin{equation}\label{eq:slackdrop}
                s_k-s_{k+1}=\frac{2c^2}{k^2},
\end{equation}
and
\begin{equation}\label{eq:constant}
 s_1=\frac{\pi^2}{48}<\sqrt2-1,
 \qquad (1+s_1)^2<2.
\end{equation}
For \(n\in B_k\), put
\begin{equation}\label{eq:construction}
 r_n=n(1+s_k),
 \qquad
 v_n=r_n\uu{\alpha_k}.
\end{equation}

We verify the restricted Fekete inequalities.  Take
\(n\in B_k\) and \(n\leq m\leq f(n)\).  By
\eqref{eq:locality}, there are only two cases.

If \(m\in B_k\), the vectors \(v_n,v_m\) point in the same direction.  If
\(n+m\in B_h\), then \(h\geq k\), so \(s_h\leq s_k\), and therefore
\[
 r_{n+m}=(n+m)(1+s_h)\leq(n+m)(1+s_k)=r_n+r_m
          =\|v_n+v_m\|_2.
\]

Suppose instead that \(m\in B_{k+1}\).  Now \(n+m\in B_h\) for some
\(h\geq k+1\), so \(s_h\leq s_{k+1}\).  Consequently, with
\(D=r_n+r_m-r_{n+m}\), equations \eqref{eq:slackdrop} and
\eqref{eq:construction} give
\begin{equation}\label{eq:radial}
 D\geq n(s_k-s_{k+1})=\frac{2c^2n}{k^2}.
\end{equation}
Hence
\begin{equation}\label{eq:radialbudget}
 (r_n+r_m)^2-r_{n+m}^2
 =D(r_n+r_m+r_{n+m})
 \geq\frac{2c^2nm}{k^2},
\end{equation}
because \(r_n+r_m+r_{n+m}\geq r_m\geq m\).

The angle between \(v_n\) and \(v_m\) is \(c/k\).  Using
\(2(1-\cos x)\leq x^2\), \(s_k,s_{k+1}\leq s_1\), and
\eqref{eq:constant}, the angular loss is bounded by
\begin{align}
 2r_nr_m\left(1-\cos\frac ck\right)
 &\leq r_nr_m\frac{c^2}{k^2}\notag\\
 &\leq (1+s_1)^2nm\frac{c^2}{k^2}
 <\frac{2c^2nm}{k^2}.
 \label{eq:angular}
\end{align}
Combining \eqref{eq:radialbudget} and \eqref{eq:angular},
\begin{align*}
 \|v_n+v_m\|_2^2-r_{n+m}^2
 &=(r_n+r_m)^2-r_{n+m}^2
   -2r_nr_m\left(1-\cos\frac ck\right)\\
 &\geq0.
\end{align*}
Thus every required inequality holds.

It remains to prove nonconvergence.  Along the block beginnings,
\[
 \frac{v_{N_k}}{N_k}=(1+s_k)\uu{\alpha_k},
 \qquad s_k\longrightarrow0,
 \qquad \alpha_k=c\sum_{j<k}\frac1j.
\]
The angles tend to infinity, while their increments tend to zero.  More
explicitly, for each \(k\) choose \(\ell(k)>k\) minimal with
\(\alpha_{\ell(k)}-\alpha_k\geq\pi\).  The divergent harmonic series makes
this possible, and minimality gives
\[
 0\leq\alpha_{\ell(k)}-\alpha_k-\pi
 <\frac{c}{\ell(k)-1}\longrightarrow0.
\]
Therefore the unit directions at indices \(k\) and \(\ell(k)\) become
asymptotically antipodal.  The subsequence \(v_{N_k}/N_k\) is not Cauchy, so
\(v_n/n\) does not converge.  This proves Theorem \ref{thm:main}. \qed

\section*{Scope, Novelty Check, and Review Focus}

The theorem is stronger than a counterexample to one proposed growth rate:
the block boundaries are constructed after seeing an arbitrary function
\(f\).  It also uses the two-dimensional Hilbert space, so no pathology of an
infinite-dimensional or merely strictly convex Banach space is involved.

A bounded novelty check through August 11, 2026 searched the exact arXiv id
and title, the wording of Question 3.2, and combinations of ``Banach-valued
Fekete lemma,'' ``window function,'' and ``uniformly convex.''  The current
arXiv record remains v1, and no later primary source resolving the question
was found.  The construction and square-summable/nonsummable device should be
checked by a specialist, with particular attention to the locality implication
\eqref{eq:locality} and the budget comparison
\eqref{eq:radialbudget}--\eqref{eq:angular}.

\begin{thebibliography}{9}
\bibitem{KS}
A.~Kulikov and F.~Shao,
\emph{Fekete's lemma in Banach spaces}, arXiv:2411.17380v1 (2024),
Question 3.2, p.~8.
\end{thebibliography}

\end{document}
